% A B S T R A C T
% ---------------

\section*{Abstract}
This thesis documents the development of \textit{Echoes of Vanitas}, a 2.5D RPG built from scratch in Godot 4 with C\#. The world is populated by anthropomorphic animals; the player assumes the role of a Bat Thief who stumbles upon a book written in the language of a long-extinct race of humans. The game was developed as a solo indie production, with all code, scenes and technical documentation produced by the author.

One pillar of the project is its original artwork. The entire visual layer, from early concept sketches to final sprites, non-player characters, inventory icons and interface elements, was hand-drawn by the author in Clip Studio Paint, with no asset packs or generators. This consistent hand-drawn 2D style is combined with the 3D lighting of the \textit{Forward Plus} renderer, producing the game's distinctive 2.5D look.

The implemented build features a surface-world hub with non-player characters, a procedurally generated dungeon (three layout variants differing in room count and lighting palette), randomised run modifiers (extra floor coins, an additional elite in the boss room, darker lighting), reward chests, four enemy types (regular foes, elites driven by a learning AI, a boss with dash and area-of-effect bombs, and invisible enemies revealed only by a radar), five mini-games in the hub (block-fit, maze, light sequence, lockpick, pipe puzzle), an inventory system with item upgrades, the sequential main quest \textit{Vanitas Cycle} (key from a puzzle cabinet, dungeon run, relic on exit, talk to the warden NPC), save/load functionality, a circular-masked minimap rendered in a \texttt{SubViewport}, and a complete user interface (HUD, tabbed menu, toasts, title screen).

The principal original contribution is \texttt{EliteBrain}, an adaptive AI module for elites modelled as a contextual multi-armed bandit. After every encounter the brain assesses how well the chosen tactic (flank, rush, kite or ambush) performed and updates its action-value table $Q$ on-line via an exponentially weighted moving average. Tactic selection follows an $\varepsilon$-greedy policy; the context is defined by two discretised variables: the player's health bucket and the number of living elites. A parallel player model (smoothed speed, back-pedal ratio, DPS) drives runtime adaptation of lead-time prediction, preferred distance and the elites' aggression and caution coefficients. The brain state is serialised to a versioned JSON file, so learning persists across sessions.

Architecturally, the project leans on a handful of classic patterns: a State Machine for player and enemy behaviour, Singletons for global services (six autoloads plus scene-bound \texttt{DungeonManager} and \texttt{UIController}), the Observer pattern realised as a static \texttt{GameEvents} class, and Godot's native component model. Dungeons are generated by a \textit{drunkard's walk} on a $9 \times 9$ grid, with special rooms (boss, treasure), unified enemy combat stats, and \texttt{NavigationAgent3D} path-baking. The in-dungeon radar uses a hemispherical \texttt{Area3D} polled every physics frame; after three seconds inside the detection volume an invisible enemy flips its shader uniform (\texttt{sprite\_visibility}) from zero to one and becomes visible.

The result is fully playable and technically complete, and it provides a solid foundation for further work, notably expanding the narrative, adding new locations, and moving from a discrete $Q$ table towards deep reinforcement learning.

\cleardoublepage
